% Приложения
\appendix

\clearpage
\phantomsection
\chapter*{ПРИЛОЖЕНИЕ А}
\addcontentsline{toc}{chapter}{ПРИЛОЖЕНИЕ А}
\renewcommand{\thelstlisting}{А.\arabic{lstlisting}}
\setcounter{lstlisting}{0}

\begin{lstlisting}[language=SQL,caption={Фрагмент реализации сущности записи на прием},label={lst:appendix-appointments}]
CREATE TABLE IF NOT EXISTS scheduling.appointments (
    id UUID PRIMARY KEY DEFAULT gen_random_uuid(),
    patient_id UUID NOT NULL REFERENCES medical.patients(id) ON DELETE RESTRICT,
    doctor_id UUID NOT NULL REFERENCES medical.doctors(id) ON DELETE RESTRICT,
    booked_service_id UUID NOT NULL REFERENCES catalog.services(id) ON DELETE RESTRICT,
    starts_at TIMESTAMPTZ NOT NULL,
    ends_at TIMESTAMPTZ NOT NULL,
    status scheduling.appointment_status NOT NULL DEFAULT 'booked',
    note TEXT,
    cancellation_reason TEXT,
    created_by_user_id UUID NOT NULL REFERENCES security.users(id) ON DELETE RESTRICT,
    cancelled_by_user_id UUID REFERENCES security.users(id) ON DELETE RESTRICT,
    completed_at TIMESTAMPTZ,
    created_at TIMESTAMPTZ NOT NULL DEFAULT NOW(),
    updated_at TIMESTAMPTZ NOT NULL DEFAULT NOW(),
    period TSTZRANGE GENERATED ALWAYS AS (tstzrange(starts_at, ends_at, '[)')) STORED,
    CHECK (ends_at > starts_at)
);
\end{lstlisting}

\begin{lstlisting}[language=SQL,caption={Фрагменты реализации ограничений целостности},label={lst:appendix-constraints}]
CREATE DOMAIN security.email_address AS VARCHAR(150)
    CHECK (VALUE IS NULL OR VALUE ~* '^[A-Z0-9._%+-]+@[A-Z0-9.-]+\.[A-Z]{2,}$');

CREATE DOMAIN security.phone_number AS VARCHAR(20)
    CHECK (VALUE IS NULL OR VALUE ~ '^\+?[1-9][0-9]{10,14}$');

ALTER TABLE catalog.service_prices
    ADD CONSTRAINT service_prices_no_overlap
    EXCLUDE USING gist (service_id WITH =, period WITH &&);

ALTER TABLE scheduling.doctor_availability
    ADD CONSTRAINT doctor_availability_no_overlap
    EXCLUDE USING gist (doctor_id WITH =, period WITH &&)
    WHERE (is_active);

ALTER TABLE scheduling.appointments
    ADD CONSTRAINT appointments_doctor_no_overlap
    EXCLUDE USING gist (doctor_id WITH =, period WITH &&)
    WHERE (status IN ('booked', 'completed'));
\end{lstlisting}

\begin{lstlisting}[language=SQL,caption={Функция получения актуальной цены услуги},label={lst:appendix-get-service-price}]
CREATE OR REPLACE FUNCTION catalog.get_service_price(
    p_service_id UUID,
    p_at TIMESTAMPTZ
)
RETURNS NUMERIC(12,2)
LANGUAGE plpgsql
AS $$
DECLARE
    v_price NUMERIC(12,2);
BEGIN
    SELECT sp.price
      INTO v_price
      FROM catalog.service_prices sp
     WHERE sp.service_id = p_service_id
       AND sp.effective_from <= p_at
       AND (sp.effective_to IS NULL OR sp.effective_to > p_at)
     ORDER BY sp.effective_from DESC
     LIMIT 1;

    IF v_price IS NULL THEN
        RAISE EXCEPTION USING ERRCODE = '22023',
            MESSAGE = format('Price is not found for service %s at %s',
                             p_service_id,
                             p_at);
    END IF;

    RETURN v_price;
END;
$$;
\end{lstlisting}

\begin{lstlisting}[language=SQL,caption={Фрагмент функции записи на прием},label={lst:appendix-book-appointment}]
CREATE OR REPLACE FUNCTION scheduling.book_appointment(
    p_patient_id UUID,
    p_doctor_id UUID,
    p_service_id UUID,
    p_starts_at TIMESTAMPTZ,
    p_created_by_user_id UUID,
    p_note TEXT DEFAULT NULL
)
RETURNS UUID
LANGUAGE plpgsql
AS $$
DECLARE
    v_duration INTEGER;
    v_ends_at TIMESTAMPTZ;
    v_appointment_id UUID;
BEGIN
    IF p_starts_at <= NOW() THEN
        RAISE EXCEPTION USING ERRCODE = '22023',
            MESSAGE = 'Appointment start time must be in the future';
    END IF;

    IF NOT EXISTS (
        SELECT 1 FROM medical.patients
         WHERE id = p_patient_id
           AND status = 'active'
    ) THEN
        RAISE EXCEPTION USING ERRCODE = '22023',
            MESSAGE = format('Patient %s is inactive', p_patient_id);
    END IF;

    IF NOT EXISTS (
        SELECT 1 FROM medical.doctors
         WHERE id = p_doctor_id
           AND status = 'active'
    ) THEN
        RAISE EXCEPTION USING ERRCODE = '22023',
            MESSAGE = format('Doctor %s is inactive', p_doctor_id);
    END IF;

    IF NOT EXISTS (
        SELECT 1
          FROM medical.doctor_services
         WHERE doctor_id = p_doctor_id
           AND service_id = p_service_id
    ) THEN
        RAISE EXCEPTION USING ERRCODE = '22023',
            MESSAGE = format('Service %s is not assigned to doctor %s',
                             p_service_id,
                             p_doctor_id);
    END IF;

    SELECT duration_minutes
      INTO v_duration
      FROM catalog.services
     WHERE id = p_service_id
       AND status = 'active';

    IF v_duration IS NULL THEN
        RAISE EXCEPTION USING ERRCODE = '22023',
            MESSAGE = format('Service %s is inactive', p_service_id);
    END IF;

    v_ends_at := p_starts_at + make_interval(mins => v_duration);

    IF NOT EXISTS (
        SELECT 1
          FROM scheduling.doctor_availability da
         WHERE da.doctor_id = p_doctor_id
           AND da.is_active
           AND da.starts_at <= p_starts_at
           AND da.ends_at >= v_ends_at
    ) THEN
        RAISE EXCEPTION USING ERRCODE = '22023',
            MESSAGE = 'Time is outside doctor availability';
    END IF;

    INSERT INTO scheduling.appointments (
        patient_id,
        doctor_id,
        booked_service_id,
        starts_at,
        ends_at,
        status,
        note,
        created_by_user_id
    )
    VALUES (
        p_patient_id,
        p_doctor_id,
        p_service_id,
        p_starts_at,
        v_ends_at,
        'booked',
        p_note,
        p_created_by_user_id
    )
    RETURNING id INTO v_appointment_id;

    RETURN v_appointment_id;
END;
$$;
\end{lstlisting}

\begin{lstlisting}[language=SQL,caption={Триггеры пересчета состояния счета},label={lst:appendix-refresh-invoice-triggers}]
CREATE OR REPLACE FUNCTION billing.trg_refresh_invoice_from_items()
RETURNS trigger
LANGUAGE plpgsql
AS $$
BEGIN
    PERFORM billing.refresh_invoice_state(COALESCE(NEW.invoice_id, OLD.invoice_id));
    RETURN COALESCE(NEW, OLD);
END;
$$;

CREATE TRIGGER trg_invoice_items_refresh_invoice
AFTER INSERT OR UPDATE OR DELETE ON billing.invoice_items
FOR EACH ROW
EXECUTE FUNCTION billing.trg_refresh_invoice_from_items();

CREATE TRIGGER trg_payments_refresh_invoice
AFTER INSERT OR UPDATE OR DELETE ON billing.payments
FOR EACH ROW
EXECUTE FUNCTION billing.trg_refresh_invoice_from_payments();
\end{lstlisting}

\begin{lstlisting}[language=SQL,caption={Реализация ролей базы данных},label={lst:appendix-db-roles}]
DO $$
BEGIN
    IF NOT EXISTS (SELECT 1 FROM pg_roles WHERE rolname = 'clinic_owner') THEN
        CREATE ROLE clinic_owner NOLOGIN;
    END IF;

    IF NOT EXISTS (SELECT 1 FROM pg_roles WHERE rolname = 'clinic_api') THEN
        CREATE ROLE clinic_api NOLOGIN;
    END IF;

    IF NOT EXISTS (SELECT 1 FROM pg_roles WHERE rolname = 'clinic_reporting') THEN
        CREATE ROLE clinic_reporting NOLOGIN;
    END IF;
END
$$;

GRANT USAGE ON SCHEMA security, catalog, medical, scheduling, billing
TO clinic_api, clinic_reporting;

GRANT SELECT, INSERT, UPDATE, DELETE ON ALL TABLES
IN SCHEMA security, catalog, medical, scheduling, billing
TO clinic_api;

GRANT EXECUTE ON ALL FUNCTIONS
IN SCHEMA security, catalog, medical, scheduling, billing
TO clinic_api;

GRANT SELECT ON ALL TABLES
IN SCHEMA security, catalog, medical, scheduling, billing
TO clinic_reporting;
\end{lstlisting}

\begin{lstlisting}[language=Go,caption={Вызов функций PostgreSQL из репозитория},label={lst:appendix-go-repository}]
func (r *AppointmentRepository) Book(
    ctx context.Context,
    command model.BookAppointmentCommand,
) (uuid.UUID, error) {
    var appointmentID uuid.UUID
    const query = "SELECT scheduling.book_appointment($1, $2, $3, " +
        "$4, $5, $6)"
    err := r.db.QueryRowContext(ctx, query,
        command.PatientID,
        command.DoctorID,
        command.ServiceID,
        command.StartsAt,
        command.CreatedByUserID,
        nullableString(command.Note),
    ).Scan(&appointmentID)

    return appointmentID, err
}

func (r *AppointmentRepository) RegisterPayment(
    ctx context.Context,
    command model.RegisterPaymentCommand,
) (uuid.UUID, error) {
    var paymentID uuid.UUID
    const query = "SELECT billing.register_payment($1, $2, $3, $4, $5)"
    err := r.db.QueryRowContext(ctx, query,
        command.InvoiceID,
        command.Amount,
        string(command.PaymentMethod),
        command.ReceivedByUserID,
        nullableString(command.ExternalReference),
    ).Scan(&paymentID)

    return paymentID, err
}
\end{lstlisting}

\begin{lstlisting}[language=SQL,caption={Проверка запрета пересечения записей},label={lst:appendix-overlap-test}]
BEGIN
    PERFORM scheduling.book_appointment(
        v_patient_2,
        v_doctor,
        v_service_consult,
        TIMESTAMPTZ '2030-01-15 10:15:00+03',
        v_patient_user,
        'SCENARIO_OVERLAP_REJECT'
    );
    v_rejected := FALSE;
EXCEPTION WHEN OTHERS THEN
    v_rejected := TRUE;
END;

IF NOT v_rejected THEN
    RAISE EXCEPTION 'Overlapping booking was not rejected';
END IF;
\end{lstlisting}

К расчётно-пояснительной записке курсовой работы прилагаются следующие материалы:
\begin{itemize}
    \item презентация к курсовой работе, состоящая из 21 слайда;
    \item USB-флеш-накопитель с исходными кодами, электронными версиями РПЗ и презентации.
\end{itemize}

\clearpage
